\section{J. 光线追踪的深层数学}

\subsection{J1. 从几何光学到波动光学}

\textbf{程函方程：}
\[
|\nabla S|^2 = n^2(\mathbf{x})
\]

\textbf{光线方程：}
\[
\frac{d}{ds}\left(n\frac{d\mathbf{x}}{ds}\right) = \nabla n
\]

\textbf{亥姆霍兹方程：}
\[
\nabla^2 \mathbf{E} + k^2 \mathbf{E} = 0, \quad k = \frac{2\pi}{\lambda}
\]

\textbf{自由空间格林函数：}
\[
G(\mathbf{x}, \mathbf{x}') = \frac{e^{ik|\mathbf{x}-\mathbf{x}'|}}{4\pi|\mathbf{x}-\mathbf{x}'|}
\]

\subsection{J2. 衍射}

\subsubsection{夫琅禾费衍射（远场）}

衍射图案是孔径函数的傅里叶变换：
\[
U(\theta_x, \theta_y) \propto \mathcal{F}\{A(x,y)\}\bigg|_{f_x = \theta_x/\lambda,\; f_y = \theta_y/\lambda}
\]

\subsubsection{菲涅尔衍射（近场）}

\[
U(x,y,z) = \frac{e^{ikz}}{i\lambda z}\iint u_0(x',y') \exp\left[\frac{ik}{2z}((x-x')^2 + (y-y')^2)\right]dx'\,dy'
\]

\subsection{J3. 偏振}

\textbf{Fresnel 方程：}
\[
r_s = \frac{n_1\cos\theta_i - n_2\cos\theta_t}{n_1\cos\theta_i + n_2\cos\theta_t}, \quad
r_p = \frac{n_2\cos\theta_i - n_1\cos\theta_t}{n_2\cos\theta_i + n_1\cos\theta_t}
\]

\textbf{Brewster 角：} $\theta_B = \arctan(n_2/n_1)$，此时 $r_p = 0$。

\subsection{J4. 相干性}

\textbf{互相干函数：}
\[
\Gamma(\mathbf{r}_1, \mathbf{r}_2, \tau) = \langle E^*(\mathbf{r}_1, t)\, E(\mathbf{r}_2, t+\tau) \rangle
\]

$|\gamma| = 1$：完全相干（激光）；$|\gamma| = 0$：完全不相干（热光源）。

大多数渲染假设不相干光（强度叠加）。

\subsection{J5. 参与介质}

\textbf{辐射传输方程（RTE）：}
\[
\frac{dL}{ds} = -\sigma_t L + \sigma_a L_e + \sigma_s \int_{S^2} p(\omega', \omega)\, L(\omega')\, d\omega'
\]

\textbf{Beer--Lambert 定律：} $T = e^{-\tau}$，$\tau = \displaystyle\int_0^s \sigma_t\, ds'$

\textbf{Delta tracking（Woodcock tracking）：} 对非均匀介质，用最大消光系数 $\sigma_{\max}$ 做包络。

\subsection{J6. 次表面散射（SSS）}

\textbf{BSSRDF：} $S(x_i, \omega_i;\, x_o, \omega_o)$，入射点和出射点可以不同。

\textbf{偶极近似：} 将次表面散射近似为两个虚拟点源，得到解析的扩散剖面。

\subsection{J7. 路径采样策略}

\begin{center}
\begin{tabular}{lll}
\toprule
\textbf{方法} & \textbf{描述} & \textbf{优势} \\
\midrule
路径追踪 (PT) & 从相机逐步采样 & 简单通用 \\
双向路径追踪 (BDPT) & 光源+相机同时构建 & 处理焦散 \\
光子映射 & 从光源发射光子 & 焦散、参与介质 \\
MLT & MCMC 在路径空间 & 复杂光路 \\
VCM/UPBP & BDPT + 光子映射统一 & 通用 \\
SPPM & 渐进收敛的光子映射 & 焦散 \\
\bottomrule
\end{tabular}
\end{center}

\subsection{J8. 焦散与灾变理论}

几何光学在焦散处预言无限亮度（振幅发散）。波动光学给出有限值（Airy 函数、Pearcey 函数）。

焦散的形态（尖点、褶皱）由灾变理论分类。